% Faithful conversion of chapter-02-text-h-0022; original XML and source characters retained.
\hypertarget{_Toc45627146}{}
\hypertarget{__RefHeading___Toc4753_3118176633}{}
\refstepcounter{section}
\section*{2.4 Conclusão}
\label{unit:chapter-02-text-h-0022}
\addcontentsline{toc}{section}{2.4 Conclusão}

% text-p-0637
\indent \originaltab{}É difícil encontrar palavras suficientes que demonstrem o quão gratificante é participar de um projeto como a AGES. A evolução observada em apenas 4 meses de envolvimento na agência é monumental. Pude trabalhar lado a lado de\originalspace{}\originalspace{}pessoas incríveis e extremamente talentosas que também estavam aprendendo e se desafiando no mesmo ambiente. Sair da minha zona de conforto, perguntar, debater, errar e refazer fizeram parte do meu dia a dia na Agência Experimental de Engenharia de Software.\originalspace{}

% text-p-0638
\indent \originaltab{}Concluo minha primeira participação na AGES com uma evolução expressiva em habilidades técnicas, interpessoais, de resoluções de conflitos e gerenciamento de tarefas. Ou seja, me torno um melhor Engenheiro de Software.

% text-p-0639
\indent \originaltab{}No começo, me senti um pouco tímido e despreparado para participar do projeto. No entanto, ao longo do desenvolvimento, ganhei mais confiança para fazer perguntas e questionamentos, e para aprender muito. Percebi também que meus colegas estavam aprendendo junto a mim, inclusive os mais experientes. O ambiente se tornou acolhedor, um espaço de aprendizado prático onde as experiências de todos se transformaram em conhecimentos para o grupo. Assim, o conhecimento de cada um pôde ser compartilhado e absorvido por todos.

% text-p-0640
\indent \originaltab{}Algumas disciplinas tiveram um papel fundamental durante minha participação na Agência Experimental de Engenharia de Software, como a disciplina de Gerenciamento de Configuração de Software, que me permitiu realizar diversas operações de versionamento de código, essenciais em um projeto dessa magnitude. Alguns conceitos que foram aprendidos de forma teórica durante meu segundo semestre no curso de Engenharia de Software foram utilizados de forma prática na AGES. Além desta disciplina, os aprendizados adquiridos nas cadeiras de Banco de Dados I e II me possibilitaram trabalhar de forma eficiente em Nosso Banco de Dados Relacional, realizando diversas operações de manipulação de dados. A cadeira de Introdução à Engenharia de Software, realizada no primeiro semestre, também teve papel crucial em nosso projeto, permitindo-me aplicar conceitos de gerenciamento de projetos na prática. Observar os AGES IV, responsável pela coordenação do time, também aplicando esses conhecimentos no ambiente prático foi extremamente enriquecedor.

% text-p-0641
\indent \originaltab{}No início do desenvolvimento do projeto, enfrentamos desafios significativos como equipe. Levou um tempo para nos ajustarmos e começarmos a trabalhar efetivamente juntos. Houve casos de tarefas realizadas de maneira incorreta, duplicação de esforços e descrições de tarefas equivocadas. No entanto, todos esses problemas foram identificados durante as \textit{Sprint Reviews }e corrigidos completamente nas Sprints subsequentes. Esse progresso reflete o comprometimento da equipe em melhorar nossa dinâmica de trabalho, tornando-a mais coesa, eficiente e agradável para todos os membros.

% text-p-0642
\indent \originaltab{}Todo o tempo e esforço dedicados pela equipe geraram resultados positivos e foram devidamente reconhecidos. Os projetos da AGES foram apresentados e avaliados por uma banca de professores, e o projeto \textit{Easy Choose} foi escolhido como Projeto Destaque. A excelente apresentação e desenvolvimento do software, especialmente pela utilização do \textit{Selenium} para realização de testes automatizados e do \textit{SonarQube} para verificação da qualidade do código, nos renderam o prêmio de melhor projeto de 2° e 4° JK em 2024/1. O registro do projeto selecionado está na Figura\hspace{0.6em}\hyperref[fig:ages-i-projeto-destaque]{\ref*{fig:ages-i-projeto-destaque}}.

% text-p-0643
\indent \originaltab{}Para finalizar, gostaria de registrar que estou muito orgulhoso da minha participação na AGES, visto que consegui realizar diversas tarefas de diferentes domínios e complexidade, desempenhando um papel efetivo e importante no meu projeto. Fiquei especialmente gratificado com a independência que adquiri ao longo do desenvolvimento. No início, necessitei de mais auxílio e orientação nas tarefas. Porém, ao longo das Sprints, ganhei muito conhecimento técnico e autonomia para conduzir minhas próprias atividades. Creio que tomarei um papel ainda mais importante e determinante em minhas próximas participações na AGES.

\memorialfigure{imagens/original/asset-0007.jpg}{\textbf{Figura 5 –} Projeto Destaque 2024/1 – AGES I}{fig:ages-i-projeto-destaque}{Fonte: pendente de identificação}{0.85}
